\begin{textsample}{NEG Dim 9 – global\_north\_2024\_09 – Score -3.00 – global\_north\_2024\_09/114507677.txt}  \label{ex:f9_neg_207}
Events this week combined to underscore my length of days, the devastation in the Holy Land -- and how the two are drawn together.
On Tuesday, my readings in the Forward Movement ministry of the Episcopal Church urged that I pray for the church in Jerusalem and the Middle East.
The next day, during a rendition of the Pledge of Allegiance at a downtown D.C. luncheon, it struck me that as a student, I recited those wordsfor 10 years before"under God"was added in 1954.
My 85th birthday on Friday marks me as older than today ' s pledge, as well as a contemporary of the travails that have plagued the Holy Land since the Arab-Israeli War of 1948.
I was a grade school student in 1950 when I learned about a Black American and U.N. official named Ralph Bunche, who won a Nobel Peace Prize for mediating armistice agreements between the new state of Israel and its four Arab neighbors.
The words"Palestine"and"Jews"might have worked their way into the conversations.
So @ @ @ @ @ @ @ @ @ @ Advertisement I was reminded of all that when I learned through this week ' s prayerful journey to the Diocese of Jerusalem that the church ' s hospital in Gaza City had been forcibly closed this summer by the Israel Defense Forces and evacuated because of escalating military action.
Crushing news.
Follow Colbert I.
King Within weeks of Hamas ' s surprise attack on Israel last year, I wrote about Al-Ahlihospital, the oldest Christian hospital in the Gaza Strip, whichtreats all people regardless of gender, race or religion.
I quoted a statement from hospital director Suhaila Tarazi to the American Friends of the Episcopal Diocese of Jerusalem that described the humanitarian crisis in Gaza as dire and changing by the hour."At this stage,"she said on Oct. 11,"our only hope is in God for a miracle in the midst of this scenery of death."Less than a year later, the church ' s mission of medical care and healing, of treating the injured and dying, was shut down. @ @ @ @ @ @ @ @ @ @ 1949 armistice, forged when I was a young boy.
The world isn't ignorant about theHamas-led attack on Israel.
We know.
We know innocent civilians were targeted, that some 1,200 were killed.
We know ofthe sexual and gender-based violence inflicted on that deadliest day in Israel ' s history -- the mutilating and despoiling of bodies, the pillaging and looting of homes, the taking of hostages.
None of that is lost on the world.
At least not the one in which I live.
Neither, however, is the scope and cost of Israel ' s vengeful response.
We know that, too."Millions of refugees remain displaced, their homes inaccessible, destroyed, or beyond repair.
Hundreds of innocents are weekly killed or severely wounded by indiscriminate attacks.
Countless others continue to endure hunger, thirst, and infectious disease.
Among these are those languishing in captivity on all sides, who additionally face the risk of ill-treatment from their captors.
Still others, far from the battlefields, @ @ @ @ @ @ @ @ @ @ and \textbf{farmlands}."Advertisement But church leaders aren't calling for retaliation.
They implore the warring factions to"reach a rapid agreement for a ceasefire resulting in the end of the war, the release of all captives, the return of the displaced, the treatment of the sick and wounded, the relief of those who hunger and thirst, and the \textbf{rebuilding} of all public and private civilian structures that have been destroyed."What of it?
Who are religious leaders to weigh in on a conflict that Israeli Prime Minister Benjamin Netanyahu regards as fundamental to his country ' s existence?
He believes that nothing less than complete and total defeat of Hamas and Hezbollah will do because they are Iran ' s two proxies to carry out the goal of Israel ' s destruction.
And so, from his viewpoint, it ' s quite all right to execute an exquisitely planned attack of \textbf{electronic} devices across Lebanon that killed at least 37 people and injured more than 3,000, even if innocent civilians end up as casualties @ @ @ @ @ @ @ @ @ @ displace the weak, hungry and innocent because they have the misfortune of living in Hamas-controlled territory.
Advertisement This column isn't about how to end the fighting or free thehostages or rebuilda war-torn land.
Frustrating and painfully slow, that is all being pursued through diplomatic channels.
Israel ' s right to exist is irrevocable.
An independent Palestinian homeland is a must.
These words, reverberating from Psalm 72, which I learned growing up, came at me this week : For he shall deliver the poor who cries out in distress, and the oppressed who has no helper.
He shall have pity on the lowly and poor ; he shall preserve the lives of the needy.
He shall redeem their lives from oppression and violence.
Something is terribly wrong.
How can the world, this nation, this over-the-hill typist, stand by when children who didn't ask to come into this world and mothers who gave birth to precious little ones are made to suffer because @ @ @ @ @ @ @ @ @ @ using the lives of others to fight to the death over... what?
More death and destruction in the name of"peace"?
And so back to that pledge of my youth : To what do we owe an allegiance?

% matched lemmas: electronic, farmland, rebuilding
\end{textsample}
